% Generated draft: requires mathematical and source review.
\documentclass[11pt]{article}
\usepackage[margin=25mm]{geometry}
\usepackage{fontspec}
\setmainfont{Latin Modern Roman}
\usepackage{amsmath,amssymb,mathtools,unicode-math}
\setmathfont{Latin Modern Math}
\usepackage{enumitem,needspace,xurl,hyperref}
\setlength{\parindent}{0pt}
\setlength{\parskip}{5pt}
\setlength{\emergencystretch}{3em}
\newcommand{\N}{\mathbb{N}}
\newcommand{\Z}{\mathbb{Z}}
\newcommand{\Q}{\mathbb{Q}}
\newcommand{\R}{\mathbb{R}}
\newcommand{\C}{\mathbb{C}}
\newcommand{\F}{\mathbb{F}}
\newcommand{\A}{\mathbb{A}}
\newcommand{\PP}{\mathbb{P}}
\newcommand{\E}{\mathbb{E}}
\newcommand{\eps}{\varepsilon}
\newcommand{\dd}{\,\mathrm{d}}
\begin{document}
\section*{UnsolvedMath ID: 30004513}
\textbf{Source problem code:} catalogue entry
\subsection{Definitions and mathematical statement}
A closure operator on \(\R^2\) is understood in the usual order-theoretic sense as a map \(\operatorname{cl}:\mathcal P(\R^2)\to\mathcal P(\R^2)\) satisfying, for all \(A,B\subseteq\R^2\),
\[
A\subseteq\operatorname{cl}(A),\qquad A\subseteq B\Longrightarrow\operatorname{cl}(A)\subseteq\operatorname{cl}(B),\qquad \operatorname{cl}(\operatorname{cl}(A))=\operatorname{cl}(A).
\]
Any further axioms intended for this operator are not specified in the supplied record. Path-connected components are taken in the subspace topology inherited from \(\R^2\). Assume that
\[
\#\{\text{path-connected components of }\operatorname{cl}(S)\}\le b
\quad\text{for every finite }S\subseteq\R^2.
\]
Here \(b\) is the component-count bound, \(K_{1,k}\) is the star with one central vertex and \(k\) leaves, and \(K_n\) is the complete graph on \(n\) vertices. The record does not explicitly specify the admissible integer ranges of \(b,k,n\) or its conventions for degenerate cases.

The principal question is whether \(kb+1\) points always suffice to obtain a constrained drawing of \(K_{1,k}\). If ``always suffice'' means a guarantee for every point set of that cardinality, the proposed assertion is: for every closure operator satisfying the displayed hypothesis and every finite \(P\subseteq\R^2\) with \(|P|=kb+1\), there exists a constrained drawing of \(K_{1,k}\) relative to \(P\) and \(\operatorname{cl}\). This is only a conditional rendering of the intended quantifiers: the supplied record does not define a constrained drawing or explain how its vertices and edges are constrained by \(P\) and \(\operatorname{cl}\). Consequently, its existence condition cannot be made into a self-contained mathematical predicate from this record alone.

Two further questions are supplied: what sharper point-count bounds hold for constrained drawings of \(K_n\), and whether constrained-drawing arguments extend to settings involving nontrivial homology or homotopy in higher dimensions. The record supplies neither a baseline bound for \(K_n\) nor the ambient spaces, dimensions, homological degrees, coefficients, homotopy groups, hypotheses, or target conclusions for the proposed extensions.

The supplied original wording refers to a closure operator on an unspecified space, whereas the clean statement specifies \(\R^2\). The planar question above follows the clean statement; the scope of any assertion for more general spaces remains unresolved.
\subsection{Short English statement}
The supplied conjecture asks whether \(kb+1\) points guarantee a constrained drawing of a star with \(k\) leaves when every finite-set closure in \(\R^2\) has at most \(b\) path-connected components. It also asks for better complete-graph bounds and extensions involving higher-dimensional homology or homotopy. The definition of a constrained drawing and the hypotheses for those extensions are missing, so this remains a qualified statement-only draft.
\subsection{Sources}
\begin{itemize}
\item Local UnsolvedMath registry, supplied record with permanent ID \(30004513\), source problem code \texttt{catalogue entry}, title ``Sharp Constrained-Drawing Bounds for Closure Operators,'' in the Oberwolfach Reports Open Problems collection. No individual proposer is identified; the supplied proposed year is 2020.
\item Supplied report citation: \emph{Discrete Geometry} (2021), \emph{Oberwolfach Reports}, pp.\ 1469--1528, DOI \texttt{10.4171/owr/2020/30}: \url{https://doi.org/10.4171/owr/2020/30}. The background describes these reports as documenting mathematical workshops and their open-problem discussions.
\item Supplied literature source: ``Oberwolfach Reports 2020/30,'' \url{https://ems.press/content/serial-article-files/46867}.
\item Literature assessment dated exactly \texttt{2026-08-22}: the supplied curation labels the problem open and states that the available constrained-star bound is quadratic in \(b\), while \(kb+1\) remains a conjectured improvement. This is supplied dated curation, not an independently verified literature assessment. The record itself requests independent verification of time-sensitive status and cited claims.
\item The registry assigns the statement status \texttt{corrected\_verified} and reports that the official report was checked, the question isolated from surrounding exposition, and its notation repaired. These are supplied provenance claims; this draft does not independently verify them or claim to have consulted either linked source.
\end{itemize}
\end{document}
